\documentclass[11pt]{article}
\usepackage{amsmath,amssymb,amsthm}
\usepackage[letterpaper,margin=1in]{geometry}
\usepackage{booktabs}
\usepackage{hyperref}
\usepackage{xcolor}
\usepackage{listings}
\usepackage{enumitem}
\usepackage{multicol}
\usepackage{float}

\definecolor{codeblue}{RGB}{40,100,180}
\definecolor{codegray}{RGB}{100,100,100}
\definecolor{codegreen}{RGB}{50,140,60}
\definecolor{backgray}{RGB}{248,248,248}

\lstset{
  language=C++,
  basicstyle=\ttfamily\small,
  keywordstyle=\color{codeblue}\bfseries,
  commentstyle=\color{codegray}\itshape,
  stringstyle=\color{codegreen},
  backgroundcolor=\color{backgray},
  frame=single, framerule=0.4pt, rulecolor=\color{codegray},
  breaklines=true, columns=fullflexible, keepspaces=true,
  showstringspaces=false, tabsize=4,
  xleftmargin=1.5em, framexleftmargin=1em,
  numbers=left, numberstyle=\tiny\color{codegray}, numbersep=8pt,
  morekeywords={constexpr,size_t,uint64_t,auto,inline,struct,concept,requires,
                expected,span,coroutine_handle,generator,mdspan,
                contract_assert,pre,post},
}

\newcommand{\var}[1]{\ensuremath{\mathrm{#1}}}
\newcommand{\svar}[1]{\ensuremath{S_{\mathrm{#1}}}}
\newcommand{\cvar}[1]{\ensuremath{C_{\mathrm{#1}}}}

\title{\textbf{Version 4 Beta --- Transition Architecture} \\[6pt]
\Large Extreme Scalability, $O(xN)$ Convergence, and Hardware Convergence \\[6pt]
\large VSEPR-SIM Development Day \#47 \\[4pt]
\normalsize From Version 3.0.X Stabilised Kernel to Version 4 Beta Formation Engine}
\author{VSEPR-SIM Development}
\date{Version 4.0.0-beta.1 --- Phase 17}

\begin{document}
\maketitle
\tableofcontents
\newpage

% ============================================================================
\section{Scope: The 3.0.X $\to$ 4.0 Transition}
\label{sec:scope}
% ============================================================================

Development Day \#47 marks the boundary between the stabilised Version~3.0.X
branch and the emerging Version~4 beta kernel.
The 3.0.X kernel can already:
\begin{itemize}[noitemsep]
  \item generate repeated structure-level outputs across multiple elemental
        systems and lattice classes (FCC, BCC, HCP),
  \item track convergence, energetics, compactness, and macro-response
        indicators in one pipeline,
  \item stream live atomistic frames via dual-port NDJSON (Phase~16),
  \item evaluate macro-property precursors (rigidity, ductility,
        thermal/electrical transport, brittleness, cohesion, fracture).
\end{itemize}

\noindent
\textbf{What 3.0.X cannot do:}
\begin{itemize}[noitemsep]
  \item distinguish a merely \emph{converged} case from a
        \emph{scientifically useful} case,
  \item quantify configurational exploration richness,
  \item rank output quality for downstream training or surrogate-model use,
  \item correlate structural compactness against macro properties,
  \item scale formation to $O(xN)$ with hardware-aware partitioning.
\end{itemize}

Version~4 beta addresses all five gaps by introducing:
\begin{enumerate}[noitemsep]
  \item the \textbf{gamma score} $\gamma$ (exploration richness),
  \item the \textbf{data quality score} $Q_\var{data}$
        (storage/publish worthiness),
  \item the \textbf{compactness score} $C_\var{compact}$
        (structural tightness),
  \item a \textbf{Pearson correlation matrix} across all final variables,
  \item an $O(xN)$ \textbf{scalable formation engine} with C++26
        contracts, erroneous-behaviour trapping, and hardware-convergent
        partitioning.
\end{enumerate}

% ============================================================================
\section{Gamma Score: Controlled Randomness and Richness}
\label{sec:gamma}
% ============================================================================

The gamma score quantifies how much meaningful configurational exploration
occurred during generation.
It does \emph{not} reward chaos; it rewards \emph{landscape coverage}.

\begin{equation}
  \label{eq:gamma}
  \boxed{
  \gamma = w_1\,\svar{bead}
         + w_2\,\svar{scale}
         + w_3\,\svar{var}
         + w_4\,\svar{intensity}
         + w_5\,\svar{orbital}
  }
\end{equation}

\subsection{Component Definitions}

\begin{description}[style=nextline,noitemsep]
  \item[$\svar{bead}$: Bead Resolution Response]
    Measures sensitivity of the final state to bead-count changes.
    For a system run at bead count $N_b$ and a reference run at $N_b'$:
    \begin{equation}
      \svar{bead} = 1 - \exp\!\Bigl(-\frac{|\eta(N_b) - \eta(N_b')|}{
                          \sigma_\eta}\Bigr)
    \end{equation}
    where $\sigma_\eta$ is the standard deviation of $\eta$ across the
    history.  A value near 0 means the result is insensitive to resolution;
    near 1 means highly resolution-dependent.

  \item[$\svar{scale}$: Multi-Scale Consistency]
    Compares derived properties at two scales (local bead level and
    macro precursor level):
    \begin{equation}
      \svar{scale} = 1 - \bigl|\rho_\var{local} - \rho_\var{macro}\bigr|
    \end{equation}
    clamped to $[0,1]$.  Rewards scale-consistency; penalises divergence
    between local packing density and macro-inferred density.

  \item[$\svar{var}$: Configurational Spread]
    The coefficient of variation of $\eta$ across all beads at the final step:
    \begin{equation}
      \svar{var} = \min\!\left(1,\;
        \frac{\sigma_\eta}{\bar{\eta}+\epsilon}\right),
      \qquad \epsilon = 10^{-12}
    \end{equation}
    Distinguishes a homogeneous crystal (low $\svar{var}$) from a
    structurally diverse morphology (high $\svar{var}$).

  \item[$\svar{intensity}$: Informational Density]
    Fraction of tracked state fields that are populated and finite:
    \begin{equation}
      \svar{intensity} = \frac{1}{F} \sum_{f=1}^{F}
        \mathbf{1}\!\bigl[\text{field}_f \neq 0 \wedge
                          \text{field}_f \in \mathbb{R}\bigr]
    \end{equation}
    where $F$ is the total number of reportable fields in the formation
    record.  A run that populates all 22 columns of the spreadsheet
    scores $\svar{intensity} = 1$.

  \item[$\svar{orbital}$: Orbital Augmentation Richness]
    Measures the active descriptor depth:
    \begin{equation}
      \svar{orbital} = \frac{\ell_\mathrm{max,active}}
                            {\ell_\mathrm{max,capable}}
    \end{equation}
    where $\ell_\mathrm{max,active}$ is the highest $\ell$ with a nonzero
    coefficient in the unified descriptor, and $\ell_\mathrm{max,capable}$
    is the system's maximum supported $\ell$.
\end{description}

\subsection{Default Weights}

\begin{center}
\begin{tabular}{ccccc}
\toprule
$w_1$ & $w_2$ & $w_3$ & $w_4$ & $w_5$ \\
\midrule
0.20 & 0.15 & 0.25 & 0.25 & 0.15 \\
\bottomrule
\end{tabular}
\end{center}

\noindent
The weights sum to 1, keeping $\gamma \in [0,1]$.

% ============================================================================
\section{Data Quality Score}
\label{sec:qdata}
% ============================================================================

The data quality score separates ``case finished'' from ``case is good.''

\begin{equation}
  \label{eq:qdata}
  \boxed{
  Q_\var{data} = a_1\,\cvar{conv}
               + a_2\,\cvar{energy}
               + a_3\,\cvar{stability}
               + a_4\,\cvar{consistency}
               + a_5\,\cvar{completeness}
               - a_6\,P_\var{failure}
  }
\end{equation}

\subsection{Component Definitions}

\begin{description}[style=nextline,noitemsep]
  \item[$\cvar{conv}$: Convergence Reward]
    Binary: 1 if the run converged within the step budget, 0 otherwise.

  \item[$\cvar{energy}$: Energy Physicality]
    Rewards negative (bound) final energy and penalises positive (unbound)
    or zero energy:
    \begin{equation}
      \cvar{energy} = \sigma\!\left(-\frac{E_\var{final}}{E_\var{ref}}\right),
      \qquad \sigma(x) = \frac{1}{1+e^{-5x}}
    \end{equation}
    where $E_\var{ref}$ is a per-element reference energy.

  \item[$\cvar{stability}$: Late-Step Stability]
    Penalises oscillation in the final 10\% of the trajectory:
    \begin{equation}
      \cvar{stability} = 1 - \min\!\left(1,\;
        \frac{\sigma_\eta^{\,\mathrm{tail}}}
             {\bar{\eta}^{\,\mathrm{tail}}+\epsilon}\right)
    \end{equation}

  \item[$\cvar{consistency}$: Cross-Field Agreement]
    Measures whether macro-precursor channels agree with structural metrics:
    \begin{equation}
      \cvar{consistency} = 1 - \frac{1}{M}\sum_{i<j}
        \bigl|\hat{x}_i - \hat{x}_j\bigr|
    \end{equation}
    for selected normalised field pairs
    $\{(\rho^*,\eta^*),\;(\var{rigidity},C^*),\;(\var{ductility},1-\var{brittleness})\}$.

  \item[$\cvar{completeness}$: Descriptor Population]
    Same as $\svar{intensity}$ from the gamma score.

  \item[$P_\var{failure}$: Failure Penalty]
    \begin{equation}
      P_\var{failure} = p_1\,\mathbf{1}[\text{step limit hit}]
                      + p_2\,\mathbf{1}[\text{NaN in output}]
                      + p_3\,\frac{n_\var{zero}}{F}
    \end{equation}
\end{description}

\subsection{Default Weights}

\begin{center}
\begin{tabular}{cccccc}
\toprule
$a_1$ & $a_2$ & $a_3$ & $a_4$ & $a_5$ & $a_6$ \\
\midrule
0.25 & 0.20 & 0.15 & 0.15 & 0.15 & 0.10 \\
\bottomrule
\end{tabular}
\end{center}

\noindent
After subtraction of the failure penalty, $Q_\var{data}$ is clamped to
$[0,1]$. A score above 0.7 qualifies a case for publication or
surrogate-model ingestion.

% ============================================================================
\section{Structure Compactness Score}
\label{sec:compact}
% ============================================================================

\begin{equation}
  \label{eq:compact}
  \boxed{
  C_\var{compact} = b_1\,\rho^*
                  + b_2\,\eta^*
                  + b_3\,R_\var{coord}^*
                  + b_4\,D_\var{macro}^*
                  - b_5\,\Phi_\var{void}^*
  }
\end{equation}

where the starred terms are normalised to $[0,1]$:

\begin{description}[style=nextline,noitemsep]
  \item[$\rho^*$: Effective Density]
    $\rho^* = \min(1,\;\bar{\rho}/\rho_\var{ref})$
    where $\rho_\var{ref}$ is the experimental crystal density.

  \item[$\eta^*$: Packing Efficiency]
    $\eta^* = \bar{\eta}$
    directly from the bead environment state (already in $[0,1]$).

  \item[$R_\var{coord}^*$: Coordination Regularity]
    $R_\var{coord}^* = 1 - \sigma_C/(\bar{C}+\epsilon)$
    Rewards uniform coordination (crystalline) and penalises irregular coordination.

  \item[$D_\var{macro}^*$: Macro Deformation Stiffness]
    $D_\var{macro}^* = (\var{rigidity\_like} + \var{cohesion\_like})/2$
    from the \texttt{MacroPrecursorState}.

  \item[$\Phi_\var{void}^*$: Void Fraction Penalty]
    $\Phi_\var{void}^* = 1 - \bar{\eta}$, the complement of packing efficiency.
\end{description}

\subsection{Default Weights}

\begin{center}
\begin{tabular}{ccccc}
\toprule
$b_1$ & $b_2$ & $b_3$ & $b_4$ & $b_5$ \\
\midrule
0.25 & 0.25 & 0.20 & 0.15 & 0.15 \\
\bottomrule
\end{tabular}
\end{center}

% ============================================================================
\section{Correlation Matrix}
\label{sec:correlation}
% ============================================================================

\subsection{Variable Set}

The clean variable set for the first-pass correlation:
\begin{equation}
  \mathbf{X} = \bigl\{
    \bar{\eta},\;\bar{\rho},\;\bar{C},\;
    \var{rms\_force},\;E_\var{final},\;
    \var{rigidity},\;\var{ductility},\;\var{cohesion},\;
    \gamma,\;Q_\var{data},\;C_\var{compact}
  \bigr\}
\end{equation}

\subsection{Pearson Correlation}

For variables $x_i$ and $x_j$ across $n$ systems:
\begin{equation}
  R_{ij} = \frac{\sum_{k=1}^{n}(x_{ik}-\bar{x}_i)(x_{jk}-\bar{x}_j)}
                {\sqrt{\sum_{k=1}^{n}(x_{ik}-\bar{x}_i)^2}
                 \sqrt{\sum_{k=1}^{n}(x_{jk}-\bar{x}_j)^2}}
\end{equation}

\subsection{Interpretation Targets}

\begin{enumerate}[noitemsep]
  \item Which variables most strongly predict compactness?
  \item Does compactness track stiffness or ductility?
  \item Does convergence quality correlate with physically meaningful outputs?
  \item Do energy and compactness move together across FCC, BCC, and HCP?
  \item Are $\gamma$ and $Q_\var{data}$ genuinely informative or decorative?
\end{enumerate}

\subsection{Reference Data (Day \#47 Spreadsheet)}

\begin{center}
\renewcommand{\arraystretch}{1.1}
\scriptsize
\begin{tabular}{llccccccccccc}
\toprule
Sym & Name & Struct & Conv & Steps & $\bar{\eta}$ & $\bar{\rho}$ &
$\bar{C}$ & $E_\mathrm{final}$ & Rigid & Duct & Coh &
$E_\mathrm{coh}$ (eV) \\
\midrule
Au & Gold       & FCC & T & 1698 & 0.279 & 5.201 & 25.66 & $-$10370 & 0.165 & 0.454 & 0.474 & $-$3.81 \\
Ag & Silver     & FCC & T &    6 & 0.001 & 0     & 0     & 0.26     & 0     & 0     & 0     & $-$2.95 \\
Cu & Copper     & FCC & T &    5 & 0.001 & 0     & 0     & 0.23     & 0     & 0     & 0     & $-$3.49 \\
Pt & Platinum   & FCC & F & 6000 & 0     & 0     & 0     & 263.3    & 0.178 & 0.433 & 0.397 & $-$5.84 \\
Ni & Nickel     & FCC & F & 5000 & 0     & 0     & 0     & 216.6    & 0.178 & 0.398 & 0.452 & $-$4.44 \\
Al & Aluminium  & FCC & T &    5 & 0.000 & 0.001 & 0     & $-$0.001 & 0.21  & 0     & 0     & $-$3.39 \\
\addlinespace
Fe & Iron       & BCC & T & 2668 & 0.469 & 11.80 & 45.34 & $-$15736 & 0.238 & 0.505 & 0.509 & $-$4.28 \\
W  & Tungsten   & BCC & T & 2954 & 0.494 & 9.394 & 40.09 & $-$33207 & 0.246 & 0.504 & 0.616 & $-$8.90 \\
Mo & Molybdenum & BCC & T & 2811 & 0.454 & 9.522 & 40.31 & $-$25013 & 0.229 & 0.495 & 0.590 & $-$6.82 \\
Cr & Chromium   & BCC & T & 2684 & 0.467 & 11.67 & 45.00 & $-$15065 & 0.237 & 0.504 & 0.607 & $-$4.10 \\
\addlinespace
Ti & Titanium   & HCP & T &   62 & 0.001 & 0.001 & 0     & $-$0.000 & 0     & 0     & 0     & $-$4.85 \\
Co & Cobalt     & HCP & F & 6000 & 0     & 0     & 0     & 253.5    & 0.091 & 0.379 & 0.361 & $-$4.39 \\
\bottomrule
\end{tabular}
\end{center}

\noindent
\textit{Observation}: BCC systems (Fe, W, Mo, Cr) converge with rich
structural content; FCC stalls (Pt, Ni) or produces near-zero descriptors
(Ag, Cu, Al); HCP is underdeveloped.
The gamma and quality scores will make these distinctions quantitative.

% ============================================================================
\section{$O(xN)$ Scalable Formation Architecture}
\label{sec:scalability}
% ============================================================================

Version~4 beta transitions the formation engine from sequential per-bead
evaluation to a partitioned parallel architecture.

\subsection{Computational Complexity}

The current kernel evaluates $N_b$ beads over $S$ steps.
Per step, each bead queries its local neighbourhood
($k$ neighbours on average):
\begin{equation}
  T_\var{step} = O(N_b \cdot k)
  \qquad\Rightarrow\qquad
  T_\var{total} = O(S \cdot N_b \cdot k)
\end{equation}

For fixed $k$ (truncated neighbour list), this is $O(S \cdot N_b)$.
The target is to reduce wall-clock time by a factor $x$ via
$x$-way partitioning:
\begin{equation}
  T_\var{parallel} = O\!\left(\frac{S \cdot N_b}{x}\right)
  + O(\log x) \;\text{(reduction)}
  = O(xN)
\end{equation}
where $x$ is the number of hardware work partitions (CPU threads or GPU SMs).

\subsection{Hardware-Convergent Partitioning}

On the target system (i9-13900K, 32 threads; RTX~4070, 46 SMs):

\begin{center}
\begin{tabular}{lll}
\toprule
\textbf{Partition} & \textbf{Hardware} & \textbf{Task} \\
\midrule
P-Cores (8)  & Raptor Cove 5.8~GHz & Formation evolution (FIRE/env update) \\
E-Cores (16) & Gracemont 4.3~GHz   & Neighbour rebuild, histogram, I/O \\
GPU SMs (46) & AD104 2.5~GHz       & Batch correlation, descriptor fitting \\
\bottomrule
\end{tabular}
\end{center}

\subsection{Domain Decomposition}

Beads are assigned to spatial cells; each cell is owned by one thread.
Ghost-zone exchange handles cross-cell neighbours:
\begin{equation}
  N_\var{cell} = \left\lceil\frac{L_\var{box}}{r_\var{cut}}\right\rceil^3,
  \qquad
  \text{beads/cell} \approx \frac{N_b}{N_\var{cell}}
\end{equation}

Communication volume scales as the surface-to-volume ratio:
\begin{equation}
  V_\var{ghost} \propto N_\var{cell}^{2/3} \cdot k
\end{equation}
which remains subdominant for $N_b > 1000$.

% ============================================================================
\section{C++26 Feature Integration}
\label{sec:cpp26}
% ============================================================================

Version~4 beta adopts the following C++26 / C++23 features:

\subsection{Runtime Contracts}

Function signatures carry executable preconditions and postconditions:
\begin{lstlisting}
auto compute_gamma(const FormationRecord& r) -> double
{
    CONTRACT_PRE(r.converged || r.steps > 0);
    double gamma = /* ... */;
    CONTRACT_POST(gamma >= 0.0 && gamma <= 1.0);
    return gamma;
}
\end{lstlisting}

\subsection{Erroneous Behaviour Trapping}

With \texttt{-ftrivial-auto-var-init=pattern}, uninitialised stack
variables receive the \texttt{0xFEFEFEFE} pattern, converting silent
data corruption into immediate, deterministic failure.

\subsection{Native Stack Traces}

Using \texttt{<stacktrace>} and \texttt{-lstdc++exp}, contract
violations produce full call-chain traces without external debuggers.

\subsection{Flat Matrix Storage}

All matrices use flat \texttt{std::vector<double>} for cache efficiency.
Row-major flat indexing: \texttt{data[i*n+j]}.

% ============================================================================
\section{Implementation Plan}
\label{sec:plan}
% ============================================================================

\begin{center}
\renewcommand{\arraystretch}{1.3}
\begin{tabular}{clll}
\toprule
\# & Task & Header & Dependencies \\
\midrule
1 & Formation record v4   & \texttt{v4/formation\_record.hpp} & Day~47 spreadsheet \\
2 & Gamma score           & \texttt{v4/gamma\_score.hpp}      & FormationRecord \\
3 & Data quality score    & \texttt{v4/data\_quality.hpp}     & FormationRecord \\
4 & Compactness score     & \texttt{v4/compactness.hpp}       & MacroPrecursorState \\
5 & Correlation matrix    & \texttt{v4/correlation\_matrix.hpp} & all three scores \\
6 & Heatmap export        & \texttt{v4/corr\_heatmap.hpp}     & CorrMatrix \\
\bottomrule
\end{tabular}
\end{center}

% ============================================================================
\section{Conclusion}
\label{sec:conclusion}
% ============================================================================

Development Day~\#47 establishes the basis for moving VSEPR-SIM from a
convergent structural kernel into a continual generation platform with
measurable informational value.
The shift from Version~3.0.X to Version~4 beta is not merely a version
increment but a methodological upgrade: each run is now judged by
convergence, variation richness, data quality, and compactness relevance.

By introducing gamma-based exploration scoring, data-quality filtering,
and a correlation matrix centred on compactness, the kernel becomes
better suited for large-scale structured generation, ranking, and
downstream thermodynamic or materials analysis.
The result is a more selective, more informative, and more scientifically
useful formation engine.

A prerequisite for reliable large-scale generation is that every
subsystem speaks the same data language. The vector type unification
completed on Development Day~\#56 (see the canonical bridge architecture
document, Section: \textit{Canonical Vector Type: \texttt{vsepr::Vec3}})
is part of this migration discipline. All layers --- atomistic kernel,
bead system, XYZ I/O, sensor, EHD, and peptide formation --- now alias
to a single authoritative \texttt{vsepr::Vec3}. This removes a class of
silent type-mismatch errors that would otherwise surface under the
increased cross-layer communication demanded by the v4 formation engine.
The Eigen adapter boundary (\texttt{to\_eigen}/\texttt{from\_eigen})
ensures that future SIMD or Eigen-backed math can be introduced without
rewriting every subsystem.

\bigskip\noindent
\textit{Humanity occasionally stumbles into competence.}

\end{document}
